\subsection{Diffusion Models as Stochastic Differential Equations}

The discrete-time formulation of denoising diffusion probabilistic models (DDPMs), proposed by
\citet{NEURIPS2020_4c5bcfec} and building on \citet{pmlr-v37-sohl-dickstein15}, defines a Markov
chain that gradually transforms data into noise through a forward noising process, and learns the
reverse transitions through variational inference. In this view, the model learns to invert a fixed
forward kernel $q(\mathbf{x}_t \mid \mathbf{x}_{t-1})$ by approximating the reverse transitions
$p_\theta(\mathbf{x}_{t-1} \mid \mathbf{x}_t)$ using the evidence lower bound (ELBO).

Around the same time, \citet{song2021scorebased} introduced \emph{score-based generative models}
(SGMs), derived from earlier work on denoising score matching \citep{vincent_score_matching_2011}.
Although SGMs and DDPMs were initially presented within different mathematical formalisms—one
discrete-time and variational, the other continuous-time and energy-based—they are now understood to
be two views of the same underlying generative mechanism. The key distinction is simply the choice
of parameterization: DDPMs work with discrete-time noise levels $\{\beta_t\}$, while SGMs treat time
as a continuous variable and parameterize corruption and denoising via a stochastic differential
equation (SDE).

\begin{figure}[H]
    \centering
    \includegraphics[width=0.65\textwidth]{chapters/background/diffusion_models/media/score.jpg}
    \caption{
        \textbf{Forward and Reverse SDEs.}
        In continuous time, the forward process that corrupts data into noise is modeled as an Itô
        SDE with drift $f(\mathbf{x},t)$ and diffusion strength $g(t)$. The reverse process, which
        recovers data from noise, has the same diffusion term but a drift modified by the
        \emph{score function} $\nabla_{\mathbf{x}} \log p_t(\mathbf{x})$, the gradient of the data
        log-density. This continuous formulation provides a unifying view of DDPMs and
        score-based models: the discrete DDPM reverse transitions approximate the numerical
        integration of this reverse SDE \citep{song2021scorebased}.
    }
    \label{fig:sde_score_diagram}
\end{figure}

\paragraph{From the Markov Chain to the Continuous Limit.}
The forward process in DDPMs is defined as a discrete Markov chain that adds Gaussian noise at each
step \eqref{eq:reparameterization_of_xt}:
\begin{equation}
\mathbf{x}_i = \sqrt{1 - \beta_i}\,\mathbf{x}_{i-1}
    + \sqrt{\beta_i}\,\boldsymbol{\epsilon}_{i-1},
\qquad
\boldsymbol{\epsilon}_{i-1} \sim \mathcal{N}(0,\mathbf{I}),
\end{equation}

To obtain a continuous formulation using a \emph{Taylor Series} expansion, we interpret the step index $i$ as discretized time
$t = i / N$ and rewrite the forward update as an infinitesimal change composed of a deterministic
component and a stochastic Brownian component \citep{song2021scorebased}. Taking the limit as $N \to \infty$ yields an Itô SDE:
\begin{equation}
\mathrm{d}\mathbf{x}
=
-\tfrac{1}{2}\beta(t)\,\mathbf{x}\,\mathrm{d}t
+
\sqrt{\beta(t)}\,\mathrm{d}\mathbf{W}(t),
\end{equation}
where $\mathbf{W}(t)$ is a Wiener process. This provides a mathematically precise continuous-time
analog of DDPM forward diffusion.

More generally, an SDE is written as
\begin{equation}
\mathrm{d}\mathbf{x}
=
f(\mathbf{x},t)\,\mathrm{d}t
+
g(t)\,\mathrm{d}\mathbf{w},
\end{equation}
with drift $f(\mathbf{x},t)$ and diffusion strength $g(t)$. The drift term describes predictable
motion (a deterministic force), while the diffusion term injects stochasticity through the Brownian
motion $\mathrm{d}\mathbf{w}$.

\paragraph{Reverse SDE and the Role of the Score.}
\citet{song2021scorebased} showed that the reverse-time dynamics of this SDE have the form
\begin{equation}
\mathrm{d}\mathbf{x}
=
\left[
    f(\mathbf{x},t)
    -
    g(t)^2 \nabla_{\mathbf{x}} \log p_t(\mathbf{x})
\right]
\mathrm{d}t
+
g(t)\,\mathrm{d}\bar{\mathbf{w}}.
\end{equation}
The term
\(
\nabla_{\mathbf{x}} \log p_t(\mathbf{x})
\)
is the \emph{score function}—the gradient of the log-density at noise level $t$. Intuitively, the
score points in the direction that makes the sample $\mathbf{x}$ more likely under $p_t$. In the
continuous formulation, this score acts as the guiding denoising vector field that pulls corrupted
samples back toward the data manifold.

This reverse SDE is the continuous-time version of the DDPM reverse Markov chain: the DDPM
mean-prediction network approximates the same underlying vector field, but expressed indirectly
through noise-prediction. Thus, introducing the SDE perspective clarifies the connection between
DDPMs and score matching, which becomes essential when extending diffusion models to tasks such as
model composition, where one must manipulate or combine score fields.

% \paragraph{Score Matching as a Theoretical Foundation.}
% Score-based generative models learn to estimate this score function directly. Classical score
% matching \citep{hyvarinen05a} aims to match the model score $s_m(\mathbf{x};\theta)$ to the data
% score $s_d(\mathbf{x}) = \nabla_{\mathbf{x}}\log p_d(\mathbf{x})$ by minimizing
% \begin{equation}
% L(\theta)
% =
% \tfrac{1}{2}
% \mathbb{E}_{p_d(\mathbf{x})}
% \!\left[
% \| s_m(\mathbf{x};\theta) - s_d(\mathbf{x}) \|^2
% \right].
% \end{equation}
% Because $s_d(\mathbf{x})$ is unknown, \citet{hyvarinen05a} derived an equivalent objective through
% integration by parts:
% \begin{equation}
% J(\theta)
% =
% \mathbb{E}_{p_d(\mathbf{x})}
% \!\left[
% \operatorname{tr}(\nabla_{\mathbf{x}} s_m(\mathbf{x};\theta))
% +
% \tfrac{1}{2}\|s_m(\mathbf{x};\theta)\|^2
% \right],
% \end{equation}
% where $\operatorname{tr}(\cdot)$ denotes the trace of the Hessian of the model score.

% Within diffusion models, this classical objective generalizes to a time-indexed version:
% the model learns the family of score functions $\nabla_{\mathbf{x}_t}\log p_t(\mathbf{x}_t)$ across
% multiple noise levels. This is precisely what the noise-prediction network in DDPM learns implicitly.

% \paragraph{Summary.}
% The SDE view provides a conceptual bridge between DDPMs and SGMs. DDPMs correspond to a specific
% discretization of the continuous forward and reverse SDEs; SGMs correspond to their continuous
% parameterization via score matching. Introducing the SDE and score viewpoint here clarifies this
% connection, which becomes important when later discussing \emph{compositional diffusion}, where the
% score plays a central role in combining or manipulating multiple generative models.
